\documentclass[10pt,a4paper]{amsart}
\usepackage{amsmath,amssymb,amsthm,graphicx}
% Retained minimal preamble for the Stacks Project chapter sites-cohomology.tex (Cohomology on Sites).
% Verbatim excerpt of the chapter preamble: Xy-pic setup (native preamble lines 14-18) and the native
% theorem declarations with the mathematical macros (native preamble lines 154-217). Copied mechanically;
% nothing added, removed or redefined. The wrapper class and the amsmath/amssymb/amsthm/graphicx base
% packages are supplied by the assembler, so the class/externaldocument part of the native preamble is omitted.
% For commutative diagrams we use Xy-pic
\usepackage[all]{xy}

% We use 2cell for 2-commutative diagrams.
\xyoption{2cell}
% Theorem environments.
%
\theoremstyle{plain}
\newtheorem{theorem}[subsection]{Theorem}
\newtheorem{proposition}[subsection]{Proposition}
\newtheorem{lemma}[subsection]{Lemma}

\theoremstyle{definition}
\newtheorem{definition}[subsection]{Definition}
\newtheorem{example}[subsection]{Example}
\newtheorem{exercise}[subsection]{Exercise}
\newtheorem{situation}[subsection]{Situation}

\theoremstyle{remark}
\newtheorem{remark}[subsection]{Remark}
\newtheorem{remarks}[subsection]{Remarks}

\numberwithin{equation}{subsection}

% Macros
%
\def\lim{\mathop{\mathrm{lim}}\nolimits}
\def\colim{\mathop{\mathrm{colim}}\nolimits}
\def\Spec{\mathop{\mathrm{Spec}}}
\def\Hom{\mathop{\mathrm{Hom}}\nolimits}
\def\Ext{\mathop{\mathrm{Ext}}\nolimits}
\def\SheafHom{\mathop{\mathcal{H}\!\mathit{om}}\nolimits}
\def\SheafExt{\mathop{\mathcal{E}\!\mathit{xt}}\nolimits}
\def\Sch{\mathit{Sch}}
\def\Mor{\mathop{\mathrm{Mor}}\nolimits}
\def\Ob{\mathop{\mathrm{Ob}}\nolimits}
\def\Sh{\mathop{\mathit{Sh}}\nolimits}
\def\NL{\mathop{N\!L}\nolimits}
\def\CH{\mathop{\mathrm{CH}}\nolimits}
\def\proetale{{pro\text{-}\acute{e}tale}}
\def\etale{{\acute{e}tale}}
\def\QCoh{\mathit{QCoh}}
\def\Ker{\mathop{\mathrm{Ker}}}
\def\Im{\mathop{\mathrm{Im}}}
\def\Coker{\mathop{\mathrm{Coker}}}
\def\Coim{\mathop{\mathrm{Coim}}}

% Boxtimes
%
\DeclareMathSymbol{\boxtimes}{\mathbin}{AMSa}{"02}

%
% Macros for moduli stacks/spaces
%
\def\QCohstack{\mathcal{QC}\!\mathit{oh}}
\def\Cohstack{\mathcal{C}\!\mathit{oh}}
\def\Spacesstack{\mathcal{S}\!\mathit{paces}}
\def\Quotfunctor{\mathrm{Quot}}
\def\Hilbfunctor{\mathrm{Hilb}}
\def\Curvesstack{\mathcal{C}\!\mathit{urves}}
\def\Polarizedstack{\mathcal{P}\!\mathit{olarized}}
\def\Complexesstack{\mathcal{C}\!\mathit{omplexes}}
% \Pic is the operator that assigns to X its picard group, usage \Pic(X)
% \Picardstack_{X/B} denotes the Picard stack of X over B
% \Picardfunctor_{X/B} denotes the Picard functor of X over B
\def\Pic{\mathop{\mathrm{Pic}}\nolimits}
\def\Picardstack{\mathcal{P}\!\mathit{ic}}
\def\Picardfunctor{\mathrm{Pic}}
\def\Deformationcategory{\mathcal{D}\!\mathit{ef}}
\usepackage{cleveref}
\crefname{originalremark}{original remark}{original remarks}
\Crefname{originalremark}{Original remark}{Original remarks}
\makeatletter
\newlabel{remark-discuss-derived-limit}{{0BKW}{source}{}{sourcepointer-1}{}}
\newlabel{remark-discuss-derived-limit@cref}{{[originalremark][0][]0BKW}{[1][0][]source}{}{}{}}
\makeatother
\crefname{originallemma}{original lemma}{original lemmas}
\Crefname{originallemma}{Original lemma}{Original lemmas}
\makeatletter
\newlabel{lemma-cohomology-over-U-trivial}{{0BKZ}{source}{}{sourcepointer-2}{}}
\newlabel{lemma-cohomology-over-U-trivial@cref}{{[originallemma][0][]0BKZ}{[1][0][]source}{}{}{}}
\makeatother
\crefname{originallemma}{original lemma}{original lemmas}
\Crefname{originallemma}{Original lemma}{Original lemmas}
\makeatletter
\newlabel{more-algebra-remark-compare-derived-limit}{{08U5}{source}{}{sourcepointer-3}{}}
\newlabel{more-algebra-remark-compare-derived-limit@cref}{{[originallemma][0][]08U5}{[1][0][]source}{}{}{}}
\makeatother
\title{Cohomology of a derived limit for a suitable inverse system}
\author{The Stacks Project Authors}
\date{stacks-project:sites-cohomology.tex:local-snapshot-20261011; GNU FDL 1.2 or later}
\begin{document}
\maketitle
\noindent\textit{Editorial note.} All delivered mathematical text is the author's own, copied from the pinned Stacks Project chapter file sites-cohomology.tex; no statement or proof step was rewritten, reordered or removed.\par\medskip
\noindent\textit{Editorial note.} The proof is the author's complete argument, the passage from the objects of B to the whole site being supplied by the covering hypothesis (1).\par\medskip
\noindent\textit{Editorial note.} Every labelled reference occurring inside the counted spans is carried in assembly.external\_references as a non-load-bearing navigation pointer (same-chapter ones at their definition site, cross-chapter ones at the invocation inside the counted span), so no load-bearing mathematics is replaced by a stub. No bibliography citation occurs inside the counted spans.\par\medskip
\noindent\textit{Editorial note.} The checked chapter list for same-name collisions is the set of stacks chapters staged in this repository: chow, cohomology, divisors, etale-cohomology, flat, more-algebra, more-morphisms, morphisms, spaces-more-morphisms, varieties together with this card's chapter sites-cohomology.\par\medskip
\noindent\textit{Editorial note.} The preamble is the chapter's own retained minimal preamble (sources/stacks-sites-cohomology/retained-preamble.tex), a verbatim excerpt of the native preamble: the Xy-pic setup and the native theorem declarations with the mathematical macros.\par\medskip
\noindent\textit{Allowed standard prerequisites (editorial).}\begin{itemize}
\item Derived limit of an inverse system of objects of the derived category. Same chapter together with the chapter Derived Categories, the defining triangle of the derived limit and its long exact cohomology sequence. Supplies the object R lim K\_n and the exact sequence used in the proof.
\item Vanishing of cohomology of the cohomology sheaves on the objects of B and vanishing of R\textasciicircum{}1 lim. Same chapter, the hypotheses (2)(a) and (2)(b) of the statement, which are the conditions under which cohomology of limits is computed by limits of cohomologies. Supplies the two vanishing statements that make the limit exact.
\end{itemize}\medskip
\ifdefined\hypertarget\hypertarget{sourcepointer-1}{}\fi
\noindent\textit{Source reference.} Original remark 0BKW: sources/stacks-sites-cohomology/pinned-original-source.tex, source lines 5201--5201. Navigation pointer invoked at proof line 5523; the counted span contains the pointer only and the cited result is not reproduced inside it, so the pointer is registered as a non-load-bearing navigation stub in the established form for Stacks material. No bibliography citation is involved.\par\medskip
\ifdefined\hypertarget\hypertarget{sourcepointer-2}{}\fi
\noindent\textit{Source reference.} Original lemma 0BKZ: sources/stacks-sites-cohomology/pinned-original-source.tex, source lines 5467--5467. Navigation pointer invoked at proof line 5524; the counted span contains the pointer only and the cited result is not reproduced inside it, so the pointer is registered as a non-load-bearing navigation stub in the established form for Stacks material. No bibliography citation is involved.\par\medskip
\ifdefined\hypertarget\hypertarget{sourcepointer-3}{}\fi
\noindent\textit{Source reference.} Original lemma 08U5: sources/stacks-sites-cohomology/pinned-original-source.tex, source lines 5532--5532. Navigation pointer invoked at proof line 5532; the counted span contains the pointer only and the cited result is not reproduced inside it, so the pointer is registered as a non-load-bearing navigation stub in the established form for Stacks material. No bibliography citation is involved.\par\medskip
\begin{lemma}
\label{lemma-derived-limit-suitable-system}
Let $(\mathcal{C}, \mathcal{O})$ be a ringed site. Let $(K_n)$
be an inverse system of objects of $D(\mathcal{O})$.
Let $\mathcal{B} \subset \Ob(\mathcal{C})$ be a subset. Assume
\begin{enumerate}
\item every object of $\mathcal{C}$ has a covering whose members are
elements of $\mathcal{B}$,
\item for all $U \in \mathcal{B}$ and all $q \in \mathbf{Z}$ we have
\begin{enumerate}
\item $H^p(U, H^q(K_n)) = 0$ for $p > 0$,
\item the inverse system $H^0(U, H^q(K_n))$ has vanishing $R^1\lim$.
\end{enumerate}
\end{enumerate}
Then $H^q(R\lim K_n) = \lim H^q(K_n)$ for $q \in \mathbf{Z}$.
\end{lemma}

\begin{proof}
Set $K = R\lim K_n$. We will use notation as in
Remark \ref{remark-discuss-derived-limit}. Let $U \in \mathcal{B}$.
By Lemma \ref{lemma-cohomology-over-U-trivial} and (2)(a)
we have $H^q(U, K_n) = H^0(U, H^q(K_n))$.
Using that the functor $R\Gamma(U, -)$ commutes with
derived limits we have
$$
H^q(U, K) = H^q(R\lim R\Gamma(U, K_n)) = \lim H^0(U, H^q(K_n))
$$
where the final equality follows from
More on Algebra, Remark \ref{more-algebra-remark-compare-derived-limit}
and assumption (2)(b). Thus $H^q(U, K)$ is the inverse limit
the sections of the sheaves $H^q(K_n)$ over $U$. Since
$\lim H^q(K_n)$ is a sheaf we find using assumption (1) that $H^q(K)$,
which is the sheafification of the presheaf $U \mapsto H^q(U, K)$,
is equal to $\lim H^q(K_n)$. This proves the lemma.
\end{proof}

\end{document}
